%Dokumentenklasse
\documentclass[
	a4paper,										%Format
	12pt,											%Schriftgröße
	twoside,										%Zweiseitig
	headsepline,								%horizontale Linie oben
	footsepline,									%horizontale Linie unten
	captions=tableheading,					%Überschrift über Tabelle (nicht unter)
]{scrartcl}

\renewcommand{\thesubsection}{\thesection.\arabic{subsection}.}
\renewcommand{\thesubsubsection}{\thesection.\arabic{subsection}.\arabic{subsubsection}. \hspace{2pt}}
\renewcommand{\theparagraph}{\thesection.\arabic{subsection}.\arabic{subsubsection}.\arabic{paragraph}.}

%Pakete:
\usepackage[exponent-product = \cdot, output-decimal-marker = {,}]{siunitx}
\usepackage{textalpha}
\usepackage{floatrow}
\floatsetup[table]{style=plaintop}
\floatsetup[longtable]{LTcapwidth=table}
%\usepackage{overcite}
\usepackage{siunitx}
\usepackage[ngerman]{babel} 			%Neue deutsche Rechtschreibung
\usepackage{grffile}
\usepackage{graphicx} 					%Bilder einbinden
\usepackage{xcolor}						%Farben für Schrift und Plots
\usepackage{pgfplots}
\pgfplotsset{compat=1.18}					%Plots und Tabellen aus Datensätzen.
\usepackage{amsmath}						%Mathe
\usepackage{amssymb}					%Mathesymbole
\usepackage{longtable}					%Seitenübergreifende Tabellen zulassen
\usepackage{float}
\usepackage{tabularx}
\usepackage{geometry}					%Seitenränder
\usepackage{ae}								%skalierbare Schriften (im PDF-File)
\usepackage[activate]{microtype}		%optischer Randausgleich
\usepackage{setspace}					%Zeilenabstand: \onehalfspacing
\usepackage[utf8]{inputenc}				%Zeichenkodierung für Umlaute
\usepackage{scrlayer-scrpage}					%für Dokumentformatierung mit Pagestyle
\usepackage[version=4]{mhchem}						%Chemische Formeln (mit \ce )
\usepackage{chemfig}
\usepackage{siunitx}
\usepackage{caption}										%Überschriften formatieren
\captionsetup{labelfont={bf,it},tablename=Tab.,figurename=Abb.}				%Überschriften formatieren
\captionsetup[table]{singlelinecheck=false,justification=raggedright}
\captionsetup[figure]{justification=raggedright}				

%Abbildungs und Tabellennummerierung anpassen
\renewcommand*\figureformat {\figurename~\arabic{section}.\arabic{subsection}.\thefigure.}	%Punkt nach Abb. #
\renewcommand*\tableformat {\tablename~\arabic{section}.\arabic{subsection}.\thetable:}		%Doppelpunkt nach Tab. #
\renewcommand*\captionformat {~}								%Nur Abstand, kein Zeichen nach Tab. #: bzw Abb. #. 

%Literatur als Subsubsection definieren und von Literaturverzeichnis in Literatur umbenennen
\makeatletter
\renewcommand*\bib@heading{\subsubsection{\bibname}			
\@mkboth{\subsubsectionmarkformat \bibname}{\subsubsectionmarkformat \bibname}}
\makeatother
\addto\captionsngerman{\renewcommand*\bibname{Literatur}} 

\usepackage{prettyref}						%Nummerierung der Tabellen und Abbildungen im Text an Tab. x.x.# anpassen
%%% Für Abbildungen %%%
\newrefformat{fig}{Abb.~\arabic{section}.\arabic{subsection}.\ref{#1}}
%%% Für Tabellen %%%
\newrefformat{tab}{Tab.~\arabic{section}.\arabic{subsection}.\ref{#1}}

%Links im Dokument:
\usepackage[									%als letztes Paktet einbinden!
	linkcolor=black,							%Farbe der Links
	citecolor=black,							%Farbe der Link-Ziele
	urlcolor=black,								%Farbe von URLs
	colorlinks=true,							%Links nicht als Kasten, sondern direkt als Text anzeigen
]{hyperref}										%Links im PDF-File, unterstützung von URLs mit /url{}

%Dokumentformatierungen
\setkomafont{sectioning}{\bfseries} 	%Serifenschrift in Überschriften
\setkomafont{subsection}{\LARGE}
\setkomafont{subsubsection}{\large}
\pagenumbering{arabic}					%arabische Seitenzahlen
\geometry{left=3cm, right=2cm, top=2.5cm, bottom=2cm} 	%Seitenränder
\parindent0pt 									%Absätze einrücken um 0Pt
\parskip 10pt									%Abstand nach Absatz
\onehalfspacing								%Zeilenabstand
\pagestyle{scrheadings}					%Kopf- und Fussleistenformatierung
\headsep0.4cm								%Abstand Kopfleiste-Text
\footskip0.8cm									%Abstand Fussleiste-Text
\urlstyle{same}								%URLs normal formatieren, nicht als Typewriter

%Befehle definieren

%%%%% create chemequation environment

\newcounter{chemequation}
\renewcommand*\thechemequation{(\Roman{chemequation})}
\newenvironment{chemequation}{%
	\stepcounter{chemequation}%
	\begin{equation}}%
{\tag*{\thechemequation}%
	\end{equation}}

%%%%% create chemtag to iterate and ref chemequations in an align environment

\newcommand*\chemtag{%
	\stepcounter{chemequation}%
	\tag*{\thechemequation}}

%%%%% to create a chemicalequation float

\usepackage{float}
\newfloat{chemicalequation}{htbp}{cex}%[chapter]
\floatname{chemicalequation}{}

\newcommand{\zitat}[1]{\textsuperscript{\cite{#1}}}

\newcommand{\bild}[4]{					%Bilder mit Beschriftung einfügen
	\begin{figure}[#4]						%#4: Platzierung (htbp)
		\centering								%Bilder mittig einfügen
		\captionbox								%Beschriftungsumgebung
			{#2\label{fig: #1}}				%#2: Beschriftung
			{\includegraphics[#3]{#1}}	%#1: Dateiname, #3: Bildeigenschaften
	\end{figure}} 

\newcommand{\formel}[2]{				%Formel einfügen und Referenz erzeugen
	\begin{equation}
		\label{eq: #1}							%#1: Formelname
		#2											%Formelinhalt
	\end{equation}}

\newcommand{\cformel}[2][]{				%Chemische Formel einügen und Referenz erzeugen
	\begin{chemequation}
		\label{ceq: #1}					%#1: Formelname (optional)
		\ce{#2}							%#2: Formelinhalt
	\end{chemequation}\\}

	% 2. Ein neuer, eigener Befehl NUR für Formeln mit Zusatz auf der rechten Seite:
\newcommand{\cformelpot}[3]{%
    \begin{chemequation}
        \label{ceq: #1}
        \ce{#2} \hspace{2cm} #3
    \end{chemequation}%
}

\newcommand{\ox}[2]{
	\overset{\text{#2}}{\ce{#1}}}
	
\newcommand{\tabelle}[6]{				%Tabelle mit Beschriftung einfügen
\begin{table}[#3]							%#3: Platzierung (htbp)
\centering										%Tabelle mittig einfügen
\captionbox										%Beschriftungsumgebung
{#2												%#2: Tabellenbeschriftung
\label{tab: #1}}								%#1: Tabellenname
{\begin{tabular}{#4}						%#4: Anzahl und Ausrichtung Spalten (l c r)
\hline 											%horizontale Linie
#5													%#5: Obere Zeile
\hline
#6													%#6: Datenzeilen
\hline 
\end{tabular}}
\end{table}}

\usepackage{array,ragged2e}
\newcolumntype{P}[1]{>{\RaggedRight\arraybackslash}p{#1}} %Verhindert ungewollte Lücken in Tabellen

\newcommand{\langetabelle}[6]{  %Seitenübergreifende Tabelle, Kopfzeile wird 
	\begin{centering}           %auf jeder neuen Seite ebenfalls angezeigt
		{\begin{longtable}[#3]{#4}
				\caption{#2\label{tab: #1}}\\				
				\hline
				#5
				\hline
				\endfirsthead
                \caption{#2\label{tab: #1}}\\
				\hline
				#5
				\hline
				\endhead
				#6												
				\hline 
		\end{longtable}}
\end{centering}}

%Plots erstellen&&&&&&&&

%#1: Referenz
%#2: Name
%#3: x label
%#4: y label
%#5: x-werte,y-werte (Spaltenweise)
%#6: Quelle (.txt Datei)
%#7: "Nummerierung" (manuell) <.<

\newcommand{\plot}[6]{
	\begin{figure}[H]	
		\centering
		\begin{tikzpicture}
	
		\begin{axis}[
		title style={at={(0.5,0)},
		anchor=north,
		yshift=-40},
		xlabel={#3},
		ylabel={#4},
		legend pos=south east,
		legend entries={},
		]
		\addplot[color=black,mark=*,only marks] table [#5] {#6};
		\label{plot: plot1}
		\end{axis}
		\end{tikzpicture}
		\caption{#2}
		\label{fig: #1}
	\end{figure}
}

\newcommand{\molL}[1]{\ensuremath{\mathrm{\SI{#1}{\frac{mol}{L}}}}} % \molL{1} => \SI{1}{\frac{mol}{L}}
\newcommand{\mL}[1]{\ensuremath{\mathrm{\SI{#1}{mL}}}} % \mL{1} => \SI{1}{mL}

\newcommand{\Rn}[1]{\romannumeral#1\relax}                   %\Rn{1} = i
\newcommand{\RN}[1]{\uppercase\expandafter{\romannumeral#1}} %\RN{1} = I

%Dokument:%%%%%%%%%%%%%%%%%%%%%%%%%%%%%%%%%%%%%%%%%%%%%%%%%%%%%%%%%%%%%%%%%%%%%%%%%%%%
\begin{document}
	
	%Titelseite
	\begin{titlepage}
		\begin{center}
			\hrule													%waagerechte Linie
			\vspace*{5cm}
			\Huge {\textbf {Versuchsreihe 5}}\\						%Titel Versuchsreihe
			\vspace*{0.5cm}
			%		\large {\bf {Untertitel}}\\						%Untertitel
			\vspace*{1cm}
			\normalsize
			Protokoll	\\									
			Praktikum für Nanowissenschaftler:innen\\
			im Studiengang Nanowissenschaften\\
			\vspace*{0.5cm}
			\Large {\textbf{Elif Öner, Friederike Penk, Maxim Popow, Michael Quint}}		\\						%Name
			\vspace*{0.5cm}
					\normalsize{Plätze 7b, 8a, 8b, 9a}								%Platznummer
			\vfill
		\end{center}
		\begin{tabbing}
			\hspace{3cm}\=\kill
			Assistent: \> Dr. Hauke Heller \\						% Name Assistent
			Version: \> 1 \\									%Version
			Datum: \> 01.11.2026 \\								%Datum
		\end{tabbing}
		\hrule														%Waagerechte Linie
\end{titlepage}

\setcounter{page}{1}
\setcounter{section}{5}								%Versuchstag #.x.x bzw Versuchsreihe
\setcounter{subsection}{0}							%Versuchsnummer x.#.x startet bei 0 für 1 
\setcounter{secnumdepth}{4}
	
\subsection{Galvanisches Element}

\subsubsection{Einleitung und Aufgabenstellung}
Ein Galvanisches Element war mit einem U-Rohr, verschiedenen Elektrolyten und Elektroden aufzubauen und die Potentialdifferenz sollte gemessen werden.
Es war der Name der aufgebauten Zelle zu benennen und der theoretische Wert mit dem Messergebnis zu vergleichen und zu diskutieren.

\subsubsection{Theoretischer Hintergrund}

	Die Grundlage für das Verständnis von galvanischen Zellen besteht darin, das 'Bestreben' von Stoffen bei Elektronenübertragungsreaktionen (auch Redox-Reaktionen genannt) zu beschreiben.
	Es werden standardmäßig Halbreaktionen als Reduktionsreaktionen geschrieben.
	\cformel[r1]{Ox + $z$e^- <=> Red}
	Dabei ist es wichtig anzumerken, dass in einer Lösung keine freien Elektronen als 'klassische' Reaktanden vorkommen. Diese Form dient als Darstellung für die in der Reaktion übertargenen Elektronen.\\
	Formal lässt sich dann der Reaktionskoeffizient $Q$ wie folgt definieren:
	\formel{Q1}{Q = \frac{a(\mathrm{Red})}{a(\mathrm{Ox})}}

	Demnach gilt für die freie Reaktionsenthalpie $\Delta_{\mathrm{R}}G$ folgende Gleichung\zitat{deltaG}:
	\formel{G1}{\Delta_{\mathrm{R}}G = \Delta_{\mathrm{R}}G^{\circ} + RT \cdot \ln Q = \Delta_{\mathrm{R}}G^{\circ} + RT \cdot \ln(\frac{a(\mathrm{Red})}{a(\mathrm{Ox})})}
	Thermodynamisch ist die gesamte umgesetzte elektrische Arbeit $W_{\mathrm{el}}$ gleich der negativen freien Reaktionsenthalpie $\Delta_{\mathrm{R}}G$;

	dabei gilt, dass die elektrische Arbeit $W_{\mathrm{el}}$ das Produkt aus der Gesamtladung $q$ und der Spannung $\Delta E$ ist, je nach Konvention des Vorzeichens:
	\formel{Arbeit}{\Delta_{\mathrm{R}}G = -W_{\mathrm{el}} = -q \cdot \Delta E}

	Dabei hängt die Gesamtladung $q$ von der Anzahl der übertragenen Elektronen pro\\ Formelumsatz $z$ und der Faraday-Konstanten $F$ (übertragene Ladung je ein Mol Elektronen) ab.
	\formel{q}{q = z\cdot F}
	\formel{Arbeit2}{\Delta_{\mathrm{R}}G = -z \cdot F \cdot \Delta E}

	Eingesetzt in die Gleichung \prettyref{eq: G1} ergibt sich dann die Nernst'sche Gleichung:
	\formel{E1}{-z \cdot F \cdot \Delta E = -z \cdot F \cdot \Delta E^{\circ} + RT \cdot \ln Q}
	\formel{E2}{\Delta E = \Delta E^{\circ} - \frac{RT}{zF} \cdot \ln Q}

	Die el. Spannung ist per Definition eine Potentialdifferenz. Das bedeutet, dass für den Vergleich immer zwei Potentiale benötigt werden. 
	Für das Standard-Reduktionspotential $E^{\circ}$ wird deshalb mittels der Standardwasserstoffelektrode (SHE) ein willkürlich gewählter Nullpunkt definiert.
	Für die Reduktion von Wasserstoff gilt demnach\zitat{H2}:
	\cformelpot{H2}{2H+ + 2e- -> H2}{E^{\circ} = 0\ \mathrm{V}}
	Für Kupfer (Cu) und Zink (Zn) gelten folgende Reduktionspotentiale\zitat{DEQ}:
	\cformelpot{Cu}{Cu^{2+} + 2e- -> Cu}{E^{\circ} = +0,34\ \mathrm{V}}
	\cformelpot{Zn}{Zn^{2+} + 2e- -> Zn}{E^{\circ} = -0,76\ \mathrm{V}}
	Da Kupfer ein positives Reduktionspotential aufweist, liegt es thermodynamisch (relativ zur Wasserstoffreduktion!) eher in reduzierter Form vor.
	Zink liegt dementsprechend eher in oxidierter Form vor.\\
	Das galvanische Element nutzt diese Reduktionspotentiale aus, um chemische Energie in elektrische Energie umzuwandeln. Dabei werden zwei sogenannte Halbzellen über einen elektrischen Leiter miteinander verbunden.
	Eine Halbzelle besteht aus einer festen Elektrode, welche in Kontakt zu einer Elektrolytlösung steht.
	Feste Metallatome gehen dann oxidiert in Lösung oder Metallionen werden reduziert und setzen sich als Feststoff auf der Elektrode ab.
	Zwei Halbzellen werden dann über eine poröse Trennwand bzw. ein Diaphragma miteinander verbunden, sodass zum Ladungsausgleich (negative) Ionen von einer Halbzelle in die andere wandern können.

	In einem Daniell-Element kommen Zink- und Kupferelektroden zum Einsatz, während eine Zinksulfat- (\ce{ZnSO4}) und eine Kupfersulfatlösung (\ce{CuSO4}) als Elektrolyt fungieren\zitat{DEEE}.
	Verbindet man nun beide Halbzellen mit einem elektrischen Leiter, finden die Reaktionen zwar räumlich getrennt voneinander statt,
	allerdings lässt sich die Gesamtreaktion als eine Redox-Reaktion aus den Gleichungen \prettyref{ceq: Cu} und \prettyref{ceq: Zn} zusammenfassen:
	\cformel[CuZn]{Cu^{2+} + 2e- + Zn -> Cu + Zn^{2+} + 2e-}
	\cformel[CuZnR]{Cu^{2+} + Zn -> Cu + Zn^{2+}}
	\bild
		{DEBild.png}					%Dateiname des Bildes, eindeutiger Name des Bildes, Verweis: \prettyref{fig: Bildname}
		{Schematische Darstellung eines Daniell-Elements: \ce{Zn|ZnSO4\ ||\ CuSO4|Cu}\\Dabei ist die Cu-Elektrode die Kathode und die Zn-Elektrode die Anode.%Text, der unter dem Bild erscheint
		\zitat{Bild}}			%Name der Quelle im Literaturverzeichnis
		{scale=1}		%angle, draft, scale, height, width
		{h}				%h: here, t: top, b: bottom, p: page; bestimmt, wo das Bild eingefügt wird. p: spezielle Seite für Abbildungen und Tabellen, Kombinationen möglich, Versuch der Platzierung gemäß Reihenfolge.

	Aus der Nernst'schen Gleichung \prettyref{eq: E2} lassen sich dann zwei Halbzellenpotentiale ermitteln:
	\formel{Cu1}{E_{\mathrm{Cu}} = E^{\circ}_{\mathrm{Cu}} - \frac{RT}{z_{\mathrm{Cu}}F} \cdot \ln Q_{\mathrm{Cu}}}
	\formel{Zn1}{E_{\mathrm{Zn}} = E^{\circ}_{\mathrm{Zn}} - \frac{RT}{z_{\mathrm{Zn}}F} \cdot \ln Q_{\mathrm{Zn}}}

	Da für beide Halbzellreakionen zwei Elektronen übertragen werden ($z=2$) und die Zellspannung die Potentialdifferenz beider Halbzellen ist, lässt sich die Zellspannung wie folgt ermitteln:
	\formel{E3}{\Delta E = E_{\mathrm{Kathode}} - E_{\mathrm{Anode}} = E_{\mathrm{Cu}} - E_{\mathrm{Zn}}}
	\formel{E4}{\Delta E = E^{\circ}_{\mathrm{Cu}} - \frac{RT}{z_{\mathrm{Cu}}F} \cdot \ln Q_{\mathrm{Cu}} - (E^{\circ}_{\mathrm{Zn}} - \frac{RT}{z_{\mathrm{Zn}}F} \cdot \ln Q_{\mathrm{Zn}})}
	\formel{E5}{\Delta E = E^{\circ}_{\mathrm{Cu}} - E^{\circ}_{\mathrm{Zn}} - \frac{RT}{2F} \cdot \ln Q_{\mathrm{Cu}} + \frac{RT}{2F} \cdot \ln Q_{\mathrm{Zn}}}
	\formel{E6}{\Delta E = \Delta E^{\circ} - \frac{RT}{2F} \cdot (\ln Q_{\mathrm{Cu}} - \ln Q_{\mathrm{Zn}})}
	\formel{E7}{\Delta E = \Delta E^{\circ} - \frac{RT}{2F} \cdot \ln (\frac{Q_{\mathrm{Cu}}}{Q_{\mathrm{Zn}}})}
	\formel{E8}{\Delta E = \Delta E^{\circ} - \frac{RT}{2F} \cdot \ln (\frac{a(\mathrm{Cu}) \cdot a(\mathrm{Zn^{2+}})}{a(\mathrm{Zn}) \cdot a(\mathrm{Cu^{2+}})})}%Bisschen unübersichtlich ich weiß :(

	Da die Aktivität von Feststoffen (Cu, Zn) per Definition 1 beträgt und man die Aktivität von Kupfer- und Zinkionen der Einfachheit halber mit deren Konzentrationen approximieren kann,
	lautet die letztendliche Gleichung:
	\formel{E9}{\Delta E = \Delta E^{\circ} - \frac{RT}{2F} \cdot \ln (\frac{[\mathrm{Zn^{2+}}]}{[\mathrm{Cu^{2+}}]})}

	Wenn Kupfersulfat und Zinksulfat in der gleichen Konzentration vorliegen, verschwindet der logarithmische Term.
	\formel{E10}{\Delta E = \Delta E^{\circ} - \frac{RT}{2F} \cdot \ln (1) = \Delta E^{\circ} - \frac{RT}{2F} \cdot 0}
	\formel{E11}{\Delta E = \Delta E^{\circ} = 0,34\ \mathrm{V} - (-0,76\ \mathrm{V}) = 1,10\ \mathrm{V}}

\subsubsection{Experimenteller Teil}

	Ein U-Rohr mit Fritte wurde auf der einen Seite mit einer blauen Kupfersulfatlösung ($1\ \mathrm{\frac{mol}{L}}$) und auf der anderen Seite mit einer farblosen Zinksulfatlösung ($1\ \mathrm{\frac{mol}{L}}$) befüllt.
	Die Kupferelektrode wurde auf der Seite der Kupfersulfatlösung fixiert und die Zinkelektrode auf der Seite der Zinksulfatlösung.
	Ein Spannungsmessgerät wurde angeschlossen und es wurde eine Gleichspannung von 1,084 V gemessen.

	\bild
		{URohr.png}					%Dateiname des Bildes, eindeutiger Name des Bildes, Verweis: \prettyref{fig: Bildname}
		{Versuchsaufbau und Messung des Daniell-Elements.%Text, der unter dem Bild erscheint
		}			%Name der Quelle im Literaturverzeichnis
		{scale=0.29}		%angle, draft, scale, height, width
		{h}				%h: here, t: top, b: bottom, p: page; bestimmt, wo das Bild eingefügt wird. p: spezielle Seite für Abbildungen und Tabellen, Kombinationen möglich, Versuch der Platzierung gemäß Reihenfolge.
\subsection{Ergebnisse und Diskussion}
\formel{Abweichung}{1 - \frac{E_{\mathrm{real}}}{E_{\mathrm{Nernst}}} = 1-\frac{1,084\ \mathrm{V}}{1,10\ \mathrm{V}}\approx 1,45\ \%}
Der abgelesene Wert aus dem Voltmeter fiel um 0,016 V geringer aus als der theoretisch ermittelte Wert, was einer Abweichung von ca. 1,45\ \% entspricht.
Mögliche Gründe sind ein Messfehler des Voltmeters, ungenaue Konzentrationen der Lösungen oder Diffusionseffekte, durch welche Kupfer- (\ce{Cu^{2+}}) bzw. Zinkionen (\ce{Zn^{2+}}) hätten diffundieren können,
wodurch die effektive Zellspannung verringert werden könnte. Außerdem ist die Verwendung der Konzentrationen in der Nernst-Gleichung eine Approximation,
da bei \ce{Cu^{2+}}- und \ce{Zn^{2+}}-Ionen selbst bei identischen Konzentrationen unterschiedliche Aktivitätskoeffizienten $\gamma_{\pm}$ vorliegen können.

\subsubsection{Zusammenfassung}
Mit Hilfe der Nernst'schen Gleichung und den Reduktionspotentialen von Kupfer und Zink ließ sich die Zellspannung des Daniell-Elements ungefähr vorhersagen.
Die Zellspannung der Zelle konnte gemessen werden und wich um ca. 1,45\ \% vom theoretischen Wert ab, was sich durch diverse kleine äußere Fehler erklären lässt.

\subsubsection{Sicherheit und Entsorgung}
		
Die für den Versuch verwendeten Chemikalien sind in \prettyref{tab: Chemikalien} mit Gefahrensymbolen, H- und P-Sätzen sowie Hinweisen zur Entsorgung aufgeführt.\zitat{test} Verwendete KMR-Stoffe sind in \prettyref{tab: KMR_leer} aufgeführt.\zitat{test}
	
\langetabelle
	{Chemikalien}
	{Verwendete Chemikalien.}
	{H}  % hier ggf. ein großes H eintragen, wenn die Tabelle sonst im Anhang verschwindet
	{l p{1.9cm} p{2.2cm} p{4.5cm} c}
	{Substanz & Symbol & H-Sätze & P-Sätze & Entsorgung \\}
	{ % Erster Stoff

  \ce{CuSO4}
    & \begin{tabular}[t]{l} GHS07 \\ GHS09 \end{tabular}
    & 302, 315 \newline 319, 410 
    & 264, 270, 273, \newline 301+312, 302+352, \newline 305+351+338
    & (1) \\

	\ce{ZnSO4.7H2O}
		& \begin{tabular}[t]{l} GHS05 \\ GHS07 \\ GHS09 \end{tabular}
		& 302, 318 \newline 410 
		&  264, 273, 280, \newline 301+312, \newline 305+351+338, 391
		& (1) \\
	}

\langetabelle
	{KMR_leer}
	{Verwendete kanzerogene, mutagene oder reproduktionstoxische Gefahrstoffe (KMR).}
	{H}
	{l P{3.8cm} P{2.4cm} P{3.5cm} c}
	{Substanz & H-Sätze & KMR-Einstufung & Verwendete Menge & Entsorgung \\}
	{ - & - & - & - & - \\}
	
Entsorgung: \\
\begin{itemize}
    \item[(1)] Die Lösung wurde im Behälter für nicht-nitrathaltige Säuren entsorgt
\end{itemize}	


\begin{thebibliography}{4} %manuelles Quellenverzeichnis

	\bibitem{deltaG}
						{E. Riedel, C. Janiak}, 
			\textit	{Anorganische Chemie}, 
						{10. Auflage}, 
						{de Gruyter}, 
						{Berlin}, 
			\textbf	{2022}, 
			S.			{315}.

	\bibitem{H2}
						{E. Riedel, C. Janiak}, 
			\textit	{Anorganische Chemie}, 
						{10. Auflage}, 
						{de Gruyter}, 
						{Berlin}, 
			\textbf	{2022}, 
			S.			{384}.

	\bibitem{DEQ}
						{E. Riedel, C. Janiak}, 
			\textit	{Anorganische Chemie}, 
						{10. Auflage}, 
						{de Gruyter}, 
						{Berlin}, 
			\textbf	{2022}, 
			S.			{379}.

	\bibitem{DEEE}
						{E. Riedel, C. Janiak}, 
			\textit	{Anorganische Chemie}, 
						{10. Auflage}, 
						{de Gruyter}, 
						{Berlin}, 
			\textbf	{2022}, 
			S.			{377}.
	
	\bibitem{Bild}
						{E. Riedel, C. Janiak}, 
			\textit	{Anorganische Chemie}, 
						{7. Auflage}, 
						{Walter de Gruyter}, 
						{Berlin}, 
			\textbf	{2007}, 
			S.			{348}.

	\bibitem{test}
						{DGUV}, 
						„{GESTIS-Stoffdatenbank}“, zu finden unter 
			\url		{https://gestis.dguv.de/}, 
			\textbf	{30.10.2026}.

			
\end{thebibliography}

%---------------------------------------------------------------------------------------------------------------------------------
\newpage

\setcounter{subsection}{1}							%Versuchsnummer x.#.x startet bei 0 für 1 
\setcounter{secnumdepth}{4}
	
\subsection{V2}

\subsubsection{Einleitung und Aufgabenstellung}


\subsubsection{Theoretischer Hintergrund}


\subsubsection{Experimenteller Teil}


\subsubsection{Ergebnisse und Diskussion}

\subsubsection{Zusammenfassung}

\subsubsection{Sicherheit und Entsorgung}
		
\langetabelle
	{Chemikalien2}
	{Verwendete Chemikalien.}
	{H}  % hier ggf. ein großes H eintragen, wenn die Tabelle sonst im Anhang verschwindet
	{l p{1.9cm} p{2.2cm} p{4.5cm} c}
	{Substanz & Symbol & H-Sätze & P-Sätze & Entsorgung \\}
	{ % Erster Stoff

    }

\langetabelle
	{KMR_leer2}
	{Verwendete kanzerogene, mutagene oder reproduktionstoxische Gefahrstoffe (KMR).}
	{H}
	{l P{3.8cm} P{2.4cm} P{3.5cm} c}
	{Substanz & H-Sätze & KMR-Einstufung & Verwendete Menge & Entsorgung \\}
	{ - & - & - & - & - \\}
	
Entsorgung: \\
\begin{itemize}
    \item[(1)]
\end{itemize}	

\begin{thebibliography}{4}		%manuelles Quellenverzeichnis

	
	\bibitem{hauke2}
						{DGUV}, 
						„{GESTIS-Stoffdatenbank}“, zu finden unter 
			\url		{https://gestis.dguv.de/}, 
			\textbf	{30.10.2026}.


\end{thebibliography}

%--------------------------------------------------------------------------------------------------------------------------------------
\newpage

\setcounter{subsection}{2}							%Versuchsnummer x.#.x startet bei 0 für 1 
\setcounter{secnumdepth}{4}
	
\subsection{V3}

\subsubsection{Einleitung und Aufgabenstellung}


\subsubsection{Theoretischer Hintergrund}


\subsubsection{Experimenteller Teil}


\subsubsection{Ergebnisse und Diskussion}

\subsubsection{Zusammenfassung}

\subsubsection{Sicherheit und Entsorgung}
		
Die für den Versuch verwendeten Chemikalien sind in \prettyref{tab: Chemikalien3} mit Gefahrensymbolen, H- und P-Sätzen sowie Hinweisen zur Entsorgung aufgeführt.\zitat{hauke3} Verwendete KMR-Stoffe sind in \prettyref{tab: KMR_leer3} aufgeführt.\zitat{hauke3}
	
\langetabelle
	{Chemikalien3}
	{Verwendete Chemikalien.}
	{H}  % hier ggf. ein großes H eintragen, wenn die Tabelle sonst im Anhang verschwindet
	{l p{1.9cm} p{2.2cm} p{4.5cm} c}
	{Substanz & Symbol & H-Sätze & P-Sätze & Entsorgung \\}
	{ % Erster Stoff

    }

\langetabelle
	{KMR_leer3}
	{Verwendete kanzerogene, mutagene oder reproduktionstoxische Gefahrstoffe (KMR).}
	{H}
	{l P{3.8cm} P{2.4cm} P{3.5cm} c}
	{Substanz & H-Sätze & KMR-Einstufung & Verwendete Menge & Entsorgung \\}
	{ - & - & - & - & - \\}
	
Entsorgung: \\
\begin{itemize}
    \item[(1)]
\end{itemize}	


\begin{thebibliography}{4}		%manuelles Quellenverzeichnis

	\bibitem{hauke3}
						{DGUV}, 
						„{GESTIS-Stoffdatenbank}“, zu finden unter 
			\url		{https://gestis.dguv.de/}, 
			\textbf	{30.10.2026}.


\end{thebibliography}

\end{document}

%\newpage

%%%%%%%%%%%%%%%%%%%%%%%%%%%%%%%%%%%%%%%%%%%%%%%%%%%%%%%%%%%%%%%
Neuer Versuch:
\subsection{Versuchsname}

\subsubsection{Einleitung und Aufgabenstellung}
\subsubsection{Theoretischer Hintergrund}
\subsubsection{Experimenteller Teil}
\subsubsection{Ergebnisse und Diskussion}
\subsubsection{Zusammenfassung}
\subsubsection{Sicherheit und Entsorgung}	
Die für den Versuch verwendeten Chemikalien sind in \prettyref{tab: entsorgung1} mit Gefahrensymbolen, H- und P-Sätzen sowie Hinweisen zur Entsorgung aufgeführt.\zitat{Gefahren1} Verwendete KMR-Stoffe sind in \prettyref{tab: KMR1} aufgeführt.
	
\langetabelle
	{entsorgung1}
	{Verwendete Chemikalien.}
	{H}  % hier ggf. ein großes H eintragen, wenn die Tabelle sonst im Anhang verschwindet
	{l p{1.4cm} p{2.2cm} p{4.5cm} c}
	{Substanz & Symbol & H-Sätze & P-Sätze & Entsorgung \\}
	{\\}
    

\langetabelle
	{KMR1}
	{Verwendete kanzerogene, mutagene oder reproduktionstoxische Gefahrstoffe (KMR).}
	{H}
	{l P{3.8cm} P{2.4cm} P{3.5cm} c}
	{Substanz & H-Sätze & KMR-Einstufung & Verwendete Menge & Entsorgung \\}
	{ - & - & - & - & - \\}
Entsorgung: \\
\begin{itemize}
    \item[(1) Die verwendete Lösung wurde mit Salzsäure angesäuert und im Behälter für nicht-nitrathaltige Säuren entsorgt.] 
\end{itemize}	

\newpage

\begin{thebibliography}{4}		%manuelles Quellenverzeichnis
	\bibitem{Gefahren1}
						{H. Heller}, 
			\textit	{Grundlagenpraktikum der Chemie für Studierende der Nanowissenschaften}, 
						{Version 1.18},
                        {Universität Hamburg},
						{Hamburg}, 
			\textbf	{2023}, 
			S.{34}.
			
			
\end{thebibliography} 
\newpage

Nützliche Beispiele:

%Tabelle für Reagenzgläserinhalte
\tabelle
	{Tabellenname}												%Eindeutige Bezeichnung der Tabelle. \prettyref{tab: Tabellenname}
		{Tabellenüberschrift.}	%Inhalt der Reagenzgläser ... bis...
		{htb}	
		{l l l}					
		{Reagenzglas & Inhalt & Beschreibung}				
		{1 & 2~mL NaOH & klare, farblose Lsg.\\
		2 & & 
		}										


%Formel Löslichkeit KClO4 -> chemische Formel
\begin{align}
    \chemtag
    \ce{KClO4(aq) <=> K+(aq) + ClO4-(aq)}
\end{align}

%hier allgemeine Formel 
\formel
		{formel2}				%eindeutige Bezeichnung der Formel, Verweis: \ref{eq: formelname}
		{L_{\mathrm{A}_{a} B_{b}} = \sqrt[m+n]{\frac{K_{\mathrm{L}}}{a^{a} \cdot b^{b}}}  }


Bilder:
	\bild
		{bildname}					%Dateiname des Bildes, eindeutiger Name des Bildes, Verweis: \prettyref{fig: Bildname}
		{Bildbeschreibung.		%Text, der unter dem Bild erscheint
		\zitat{Quelle}}			%Name der Quelle im Literaturverzeichnis
		{Bildeigenschaften}		%angle, draft, scale, height, width
		{Platzierung}				%h: here, t: top, b: bottom, p: page; bestimmt, wo das Bild eingefügt wird. p: spezielle Seite für Abbildungen und Tabellen, Kombinationen möglich, Versuch der Platzierung gemäß Reihenfolge.
	
		
	Bildeigenschaften:
		angle - Bestimmt den Drehwinkel um einen Referenzpunkt für die Grafik. Positive Werte drehen die Grafik gegen den Uhrzeigersinn, negative mit ihm. 
		draft - Ist draft gesetzt, wird nur ein Platzhalter der Grafik angezeigt. Beschleunigt den Übersetzungsvorgang des Dokuments erheblich. 
		scale - Skaliert die Grafik anhand eines Skalierungsfaktors. 
		height - Skaliert die Grafik auf die angegebene Höhe. Die Angabe der Größe muss in einer LaTeX-spezifischen Einheit erfolgen. 
		width - Skaliert die Grafik auf die angegebene Breite. Die Angabe der Größe muss in einer LaTeX-spezifischen Einheit erfolgen. 	
	

Formeln:
	\formel
		{formelname}				%eindeutige Bezeichnung der Formel, Verweis: \ref{eq: formelname}
		{Formel}						%


Tabellen:
	\tabelle
		{Tabellenname}												%Eindeutige Bezeichnung der Tabelle. \prettyref{tab: Tabellenname}
		{Tabellenüberschrift.\zitat{Quelle}}						%Text, der über der Tabelle erscheint, Quelle der Daten
		{Platzierung}													%h: here, t: top, b: bottom, p: page; bestimmt, wo das Bild eingefügt wird. p: spezielle Seite für Abbildungen und Tabellen, Kombinationen möglich, Versuch der Platzierung gemäß Reihenfolge.		
		{Spaltenanzahl_und_Ausrichtung}						%{ l c r }: 3 Spalten, links mittig und rechts; { c c c c }: 4 Spalten, alle mittig
		{Titelzeile}														%Obere Zeile(n), danach folgt ein waagerechter Strich, Spalten durch & trennen, Neue Zeile durch \\
		{Datenzeilen}													%Zeilen unter dem waagerechten Strich, Spalten durch & trennen, Neue Zeile durch \\
		
	Zellen verbinden:	
		\multicolumn{Anzahl_der_verbunden_Spalten}{Ausrichtung (l/c/r}{Inhalt}
	
Reaktionen:
Formeln: 
\begin{center}
\ce{NH+4}, \ce{Fe3(CO)12}, \ce{Al^{3+}},\ce{Cr2O^{2-}7}
\end{center}

Isotope:
\begin{center}
\ce{^{37}_{17}Cl}, \ce{^{32}_{16}S^{2-}}
\end{center}

Bindungen: 
\begin{center}
\ce{H-O-N=S}, \ce{Cl-C#N}
\end{center}

Reaktionsgleichungen: 
\begin{center}
\ce{2H2 + O2 -> 2H2O}
\ce{MgCl2 ->[\text{Elektrolyse}] Mg\,(l) + Cl_{2}\,(g)}
\end{center}

Redoxreaktion:
\begin{center}
$\overset{0}{\ce{S}} + \overset{0}{\ce{O2}} \ce{->} \overset{+\text{IV}}{\ce{S}}+\overset{-\text{II}}{\ce{O2}}$
\end{center}


Literaturverweise:
		\bibitem{zeitschrift}
						{A. Autor}, 
			\textit	{Zeitschrift} 
			\textbf	{Jahr}, 
			\textit	{Band} 
						{Seite-Seite}.
		
		\bibitem{buch_ohne_herausgeber}
						{A. Autor}, 
			\textit	{Titel}, 
						{Verlag}, 
						{Ort}, 
			\textbf	{Jahr}, 
			S.			{Seite-Seite}.
			
		\bibitem{buch_mit_herausgeber}
						{A. Autor}, 
			\textit	{Titel}, 
						{Auflage}, 
						{Verlag}, 
						{Ort}, 
			\textbf	{Jahr}, 
			S.			{Seite-Seite}.
			
		\bibitem{dissertation}
						{A. Autor}, 
						{Dissertation}, 
						{Universität}, 
			\textbf	{Jahr}.
			
		\bibitem{internetseite}
						{A. Autor}, 
						„{Titel}“, zu finden unter 
			\url		{www.xyz.de}, 
			\textbf	{Jahr}.

% Kommentare Sophie:

% <- Für Kommentare oder Code der gerade nicht gebraucht wird, wird beim kompilieren nicht beachtet
%Um Bilder einzufügen: \bild{'Dateiname' (Datei muss im gleichem Ordner mit der TeX Datei liegen)}{'Bildunterschrift'}{width= ...cm (Größe des Bildes)}{H (Ort des Bildes, H=Hier)}

%Für Summenformeln im Text: \ce{z.B.: H2SO4 (Zahlen werden automtisch tief gestellt)}

% Um Formeln im Text zu referieren: \prettyref{fig (für Bilder und Graphen)/ eq (für mathematische Formeln)/ ceq (für chemische Formeln): 'Referenz'}

%Für chemische Formeln: \cformel['Referenz']{2\ox{Fe}{+\RN{3}}Cl3 + \ox{Fe}{+\RN{2}}Cl2 + 8NH4+ + 8 OH- \ce{<=>[{N2H4}]} \ox{Fe}{+\RN{2}}O\cdot \ox{Fe}{+\RN{3}}2O3 \ce{v} + 4 H2O + 8NH4+ + 4 Cl-}

%Für mathematische Formeln: \formel{'Referenz'}{2\cdot d\cdot \ce{sin}(\theta)= n \lambda}

%Tabellen: Online Converter um Exel Tabellen in LaTeX Format zu bringen: https://www.tablesgenerator.com
%\tabelle{'Referenz'}{'Tabellenüberschrift'}{H}{'Breite der Spalte bzw. Textausrichtung' (Spaltenbreite passt sich Textlänge an: l=links, r=rechts, c=zentriert. Spaltenbreite fest und Blocksatz: p{...cm}), Jede Spalte muss einzeln definiert werden, z.B l l l = 3 Spalten}{'Spaltenüberschrift' muss der Anzahl an Spalten entsprechen. Spaltenüberschriften müssen durch ein & getrennt werden. Außerdem wichtig am Ende: \\}{'Tabelleneinträge', müssen ebenfalls durch ein & getrennt werden und am Ende der Zeile ein \\ haben
%	- & - & - & \\
%	- & - & - & \\}
%Beispiel für eine Tabelle (zum Ausführen '%' entfernen):
%\tabelle{Auflösungsvermögen}{Experimentell ermitteltes und theoretisches Auflösungsvermögen für Objektiv 20X und 63X bei schräger und gerader Beleuchtung.}{H}{l l p{2cm} p{3cm}  p{3cm}}{Objektiv&Beleuchtung& Spalt mit einem Peak&experimentelles Auflösungsvermögen~/~nm&theoretisches Auflösungsvermögen~/~nm\\}
%	{20X&gerade&16&>750&1125\\
%	20X&schräg&9&690&563\\
%	63X&gerade&11&>750&600\\
%	63X&schräg&7&584&300\\}